% GENERATED FILE. Edit canonical lessons, then run npm run generate:books.
% canonical-source-hash: fnv1a64:f5fcd17c84c50f36
% canonical-lessons: FA-C425-otoban, FA-C425-avarezi, FA-C425-shane-khaki, FA-C425-terafik, FA-C425-monharef-shodan

\chapter{On the motorway}
\label{ch:fa-on-the-motorway}
\hlchaptermodality{425}

\begin{chapteropening}
\textbf{By the end of this chapter:} \emph{I can talk about the motorway: the motorway itself, the toll, the hard shoulder, the traffic and swerving.}

Complete the last lesson of chapter 425: I can talk about the motorway: the motorway itself, the toll, the hard shoulder, the traffic and swerving.
\end{chapteropening}

\section[\texorpdfstring{\fa{اتوبان} (otobân)}{اتوبان (otobân)}]{\fa{اتوبان} (otobân) — a motorway (US: a freeway)}

\label{lesson:FA-C425-otoban}



\begin{quote}
Before the new one: say the Persian for to pull over, then the Persian for safety.
\end{quote}

\subsection*{You'll want to know: \fa{اتوبان}}
\textbf{\fa{اتوبان}} (\emph{otobân}) — \textquotedblleft{}a motorway (US: a freeway)\textquotedblright{}.

\begin{quote}
\fa{از} \fa{این} \fa{شهر} \fa{تا} \fa{پایتخت} \fa{اتوبان} \fa{است}.
\end{quote}

\begin{quote}
\emph{Az in shahr tâ pâytakht otobân ast.}
\end{quote}

\begin{quote}
There is a motorway from this town to the capital.
\end{quote}

The more formal word is \emph{bozorgrâh}.

\begin{cousinweb}
\emph{otobân} was borrowed from German \emph{Autobahn} 'motorway' (EtymDB), a compound of \emph{Auto} 'car' and \emph{Bahn} 'way'.
\end{cousinweb}

\begin{grammarlens}[title={What you've built}]
One more word for telling, explaining and describing at B1.
\end{grammarlens}

\subsection*{Your turn}
\begin{itemize}
\raggedright
  \item \emph{Say it:} \emph{otobân}
  \item \emph{Say it:} \emph{otobân}, once more
  \item \emph{Say it:} a sentence of your own with \emph{otobân}
  \item \emph{Recall:} say the Persian for to pull over, then the Persian for safety, then say \emph{otobân} again
\end{itemize}

\begin{tcolorbox}[breakable,colback=teal!4,colframe=teal!35!black,title={Before you move on}]
What does \fa{اتوبان} mean? (\textquotedblleft{}A motorway (US: a freeway)\textquotedblright{}.) Say it once more.
\end{tcolorbox}
\section[\texorpdfstring{\fa{عوارضی} ('avârezi)}{عوارضی ('avârezi)}]{\fa{عوارضی} ('avârezi) — a toll booth, a toll gate}

\label{lesson:FA-C425-avarezi}



\begin{quote}
Before the new one: say the Persian for safety, then the Persian for a motorway.
\end{quote}

\subsection*{You'll want to know: \fa{عوارضی}}
\textbf{\fa{عوارضی}} (\emph{'avârezi}) — \textquotedblleft{}a toll booth, a toll gate\textquotedblright{}.

\begin{quote}
\fa{در} \fa{عوارضی} \fa{ماشینها} \fa{کمی} \fa{صبر} \fa{میکنند}.
\end{quote}

\begin{quote}
\emph{Dar 'avârezi mâshin-hâ kami sabr mikonand.}
\end{quote}

\begin{quote}
At the toll gate the cars wait a little.
\end{quote}

Built on \emph{'avârez} 'tolls, charges', an Arabic broken plural, + \emph{-i}.

\begin{cousinweb}
\emph{'avârezi} is built from \emph{'avârez} 'tolls, charges', a loan usually traced to the Arabic plural \emph{ʿawāriḍ} 'things that come up, extra charges', and the Persian ending \emph{-i}. EtymDB does not list it.
\end{cousinweb}

\begin{grammarlens}[title={What you've built}]
One more word for telling, explaining and describing at B1.
\end{grammarlens}

\subsection*{Your turn}
\begin{itemize}
\raggedright
  \item \emph{Say it:} \emph{'avârezi}
  \item \emph{Say it:} \emph{'avârezi}, once more
  \item \emph{Say it:} a sentence of your own with \emph{'avârezi}
  \item \emph{Recall:} say the Persian for safety, then the Persian for a motorway, then say \emph{'avârezi} again
\end{itemize}

\begin{tcolorbox}[breakable,colback=teal!4,colframe=teal!35!black,title={Before you move on}]
What does \fa{عوارضی} mean? (\textquotedblleft{}A toll booth, a toll gate\textquotedblright{}.) Say it once more.
\end{tcolorbox}
\section[\texorpdfstring{\fa{شانه} \fa{خاکی} (shâne-khâki)}{شانه خاکی (shâne-khâki)}]{\fa{شانه} \fa{خاکی} (shâne-khâki) — a hard shoulder (of a road)}

\label{lesson:FA-C425-shane-khaki}



\begin{quote}
Before the new one: say the Persian for a motorway, then the Persian for a toll booth.
\end{quote}

\subsection*{You'll want to know: \fa{شانه} \fa{خاکی}}
\textbf{\fa{شانه} \fa{خاکی}} (\emph{shâne-khâki}) — \textquotedblleft{}a hard shoulder (of a road)\textquotedblright{}.

\begin{quote}
\fa{ماشین} \fa{خراب} \fa{را} \fa{در} \fa{شانه} \fa{خاکی} \fa{دیدیم}.
\end{quote}

\begin{quote}
\emph{Mâshin-e kharâb râ dar shâne-khâki didim.}
\end{quote}

\begin{quote}
We saw the broken-down car on the hard shoulder.
\end{quote}

Literally 'earthy shoulder': \emph{shâne} 'shoulder' + \emph{khâki} 'of earth'.

\begin{cousinweb}
\emph{shâne-khâki} is a native compound: \emph{shâne} 'shoulder' is inherited from Middle Persian \emph{šānag} (EtymDB), and \emph{khâki} 'earthy' is built on \emph{khâk} 'earth', from Middle Persian \emph{xāk} (EtymDB).
\end{cousinweb}

\begin{grammarlens}[title={What you've built}]
One more word for telling, explaining and describing at B1.
\end{grammarlens}

\subsection*{Your turn}
\begin{itemize}
\raggedright
  \item \emph{Say it:} \emph{shâne-khâki}
  \item \emph{Say it:} \emph{shâne-khâki}, once more
  \item \emph{Say it:} a sentence of your own with \emph{shâne-khâki}
  \item \emph{Recall:} say the Persian for a motorway, then the Persian for a toll booth, then say \emph{shâne-khâki} again
\end{itemize}

\begin{tcolorbox}[breakable,colback=teal!4,colframe=teal!35!black,title={Before you move on}]
What does \fa{شانه} \fa{خاکی} mean? (\textquotedblleft{}A hard shoulder (of a road)\textquotedblright{}.) Say it once more.
\end{tcolorbox}
\section[\texorpdfstring{\fa{ترافیک} (terâfik)}{ترافیک (terâfik)}]{\fa{ترافیک} (terâfik) — traffic; a traffic jam}

\label{lesson:FA-C425-terafik}



\begin{quote}
Before the new one: say the Persian for a toll booth, then the Persian for a hard shoulder.
\end{quote}

\subsection*{You'll want to know: \fa{ترافیک}}
\textbf{\fa{ترافیک}} (\emph{terâfik}) — \textquotedblleft{}traffic; a traffic jam\textquotedblright{}.

\begin{quote}
\fa{صبحها} \fa{در} \fa{اتوبان} \fa{ترافیک} \fa{سنگین} \fa{است}.
\end{quote}

\begin{quote}
\emph{Sobh-hâ dar otobân terâfik sangin ast.}
\end{quote}

\begin{quote}
In the mornings the traffic on the motorway is heavy.
\end{quote}

\emph{Dar terâfik mândan} — to be stuck in traffic.

\begin{cousinweb}
\emph{terâfik} was borrowed from a European language: it is usually traced to English \emph{traffic} or French \emph{trafic}. EtymDB does not list the Persian word.
\end{cousinweb}

\begin{grammarlens}[title={What you've built}]
One more word for telling, explaining and describing at B1.
\end{grammarlens}

\subsection*{Your turn}
\begin{itemize}
\raggedright
  \item \emph{Say it:} \emph{terâfik}
  \item \emph{Say it:} \emph{terâfik}, once more
  \item \emph{Say it:} a sentence of your own with \emph{terâfik}
  \item \emph{Recall:} say the Persian for a toll booth, then the Persian for a hard shoulder, then say \emph{terâfik} again
\end{itemize}

\begin{tcolorbox}[breakable,colback=teal!4,colframe=teal!35!black,title={Before you move on}]
What does \fa{ترافیک} mean? (\textquotedblleft{}Traffic; a traffic jam\textquotedblright{}.) Say it once more.
\end{tcolorbox}
\section[\texorpdfstring{\fa{منحرف} \fa{شدن} (monharef shodan)}{منحرف شدن (monharef shodan)}]{\fa{منحرف} \fa{شدن} (monharef shodan) — to swerve, to veer off; to go astray}

\label{lesson:FA-C425-monharef-shodan}



\begin{quote}
Before the new one: say the Persian for a hard shoulder, then the Persian for traffic.
\end{quote}

\subsection*{You'll want to know: \fa{منحرف} \fa{شدن}}
\textbf{\fa{منحرف} \fa{شدن}} (\emph{monharef shodan}) — \textquotedblleft{}to swerve, to veer off; to go astray\textquotedblright{}.

\begin{quote}
\fa{ماشین} \fa{روی} \fa{برف} \fa{منحرف} \fa{شد}.
\end{quote}

\begin{quote}
\emph{Mâshin ru-ye barf monharef shod.}
\end{quote}

\begin{quote}
The car swerved on the snow.
\end{quote}

\emph{monharef} 'turned aside' + \emph{shodan}. Also figurative: to go off the right path.

\begin{cousinweb}
\emph{monharef shodan} is built from \emph{monharef} 'turned aside', a loan usually traced to Arabic \emph{munḥarif}, and the native verb \emph{shodan} 'to become'. EtymDB does not list it.
\end{cousinweb}

\begin{grammarlens}[title={What you've built}]
One more word for telling, explaining and describing at B1.
\end{grammarlens}

\subsection*{Your turn}
\begin{itemize}
\raggedright
  \item \emph{Say it:} \emph{monharef shodan}
  \item \emph{Say it:} \emph{monharef shodan}, once more
  \item \emph{Say it:} a sentence of your own with \emph{monharef shodan}
  \item \emph{Recall:} say the Persian for a hard shoulder, then the Persian for traffic, then say \emph{monharef shodan} again
\end{itemize}

\begin{tcolorbox}[breakable,colback=teal!4,colframe=teal!35!black,title={Before you move on}]
What does \fa{منحرف} \fa{شدن} mean? (\textquotedblleft{}To swerve, to veer off; to go astray\textquotedblright{}.) Say it once more.
\end{tcolorbox}
